\begin{table}[!htb]
\centering
\renewcommand\arraystretch{1.15}
\setlength{\tabcolsep}{0.14cm}
\caption{ULTRASAT (LCS grid) and UVEX under identical scheduling assumptions:
the same simulated events, the same kilonova integrated through each mission's
NUV band, SNR\,$=5$, two visits, observations between 15\,min and 6\,hr after
the merger. All values are expected numbers per year, $\lambda$, at the median
merger rate; the 90\% credible intervals are given in
Tables~\ref{tab:ultrasat-non-overlap-selected-detected-1yr}
and~\ref{tab:uvex-selected-detected-1yr}. \emph{GW}: BNS and NSBH detections.
\emph{Sel.}: selected events; \emph{Det.}: detected kilonovae.
\emph{Ratio}: UVEX over ULTRASAT detected counts. Because both missions follow
the same events, the merger-rate uncertainty cancels in this ratio.}
\label{tab:mission-comparison}
\begin{tabular}{lcccccc}
\toprule
 & & \multicolumn{2}{c}{ULTRASAT} & \multicolumn{2}{c}{UVEX} & \\
\cmidrule(lr){3-4} \cmidrule(lr){5-6}
Run & GW & Sel. & Det. & Sel. & Det. & Ratio \\
\midrule
O5a & 19.5 & 6.68 & 2.46 & 8.69 & 3.52 & 1.4 \\
O5b & 30.6 & 8.12 & 2.61 & 10.52 & 4.61 & 1.8 \\
O5c & 38.0 & 9.64 & 3.78 & 13.77 & 6.43 & 1.7 \\
\midrule
O5 total & 88.0 & 24.44 & 8.85 & 32.97 & 14.56 & 1.6 \\
\bottomrule
\end{tabular}
\end{table}
